% !TeX root = ../教材题库_全部条目_转写工作稿.tex
% Source: XJ-B1-C05-TH-P170-Y0981
% 第5章 / 5.2 任意角的三角函数 / 教材 p.170 / PDF 第177页
% OCR未校对：公式与文字均可能有误；未完成数学TeX转写；禁止直接发给学生
% 原题及相邻原解截图（核对用）：../题目原图/XJ-B1-C05-TH-P170-Y0981-01.webp
% 本文件只在转写完成后作为题干使用；OCR 条目尚不是可靠题干。
\begin{bankitem}{p.170 思考（含紧邻说明）}{第5章 / 5.2 任意角的三角函数 / 教材 p.170 / PDF 第177页}{XJ-B1-C05-TH-P170-Y0981}{OCR未校对：公式与文字均可能有误；未完成数学TeX转写；禁止直接发给学生}
\begin{ocrdraft}
[教材p.170 / PDF第177页]
\par
？思考在平面直角坐标系内，以原点0为圆心作一单位圆．设任意角的终
\par
边与单位圆相交于点P．
\par
（1)作点P关于轴的对称点PI,则ZxOPl与角a有什么关系？Z主0P1的三
\par
角函数值与角a的三角函数值之间有什么关系？
\par
（2）作点P关于原点的对称点P2，则Z0P2与角有什么关系？Z主0P2的三
\par
角函数值与角a的三角函数值之间有什么关系？
\par
设任意角的终边与单位圆的交点坐标为P(x，y)，则
\par
当x0时，tan仅一
\par
Sln．一3，，COS．一X，
\par
（1)如图5，2一8，作点P(x，y)关于轴的对称点
\par
00，一y)，则／鼠丿PI
\par
由三角函数的定义可得
\par
0，过）
\par
Sln，
\par
图5．2一8
\par
当x0时，
\par
tan(
\par
一一tan．
\end{ocrdraft}
\end{bankitem}
